% =====================================================================
% SUPPLEMENTARY METHODS — the method detail the main text points at.
%
% SCOPE, AW 15 November: "the Methods material Results points at goes
% here". Four items: Severity and Label transfer (Results points at both),
% the gene-group construction in full (Note 1 and Note 3 point at it), and
% the training hyperparameters with a generated table (AW's request).
% Ordered as the main Methods: data, gene vocabulary, optimisation, then
% the analysis built on the trained model.
%
% REGISTER, read before writing (17 Nov): MIDAS's Supplementary Methods
% and IBDex's Methods in full, scGPT Note 12 for definitions, and
% SUPPLEMENTARY_REGISTER.md. First person for procedure, past tense,
% procedure before outcome, every setting stated with its value, a
% one-clause reason per choice, connectives at the joints, no semicolon
% chains, no internal vocabulary.
%
% EVERY NUMBER HAS A FILE.
%   Severity        merge_columns(SEVERITY / SEVERITY_CASES) over the cell
%                   metadata: 641,024 cells, 61/29/28 donors; 333,984
%                   cells, 57 patients. Supplementary Data 3.
%   Vocabulary      D5_pathway_groups_and_genes.csv: 44 groups, 247-493,
%                   median 341.5; Reactome 28 groups 10,114 genes.
%                   build_vocabulary.sh: --cut union, 250/500.
%   Hyperparameters tables/hyperparameters.tex, generated by
%                   figures/scripts/supp_methods.py from run_config.json.
%   Label transfer  transfer_cell_family.py (K = 25); x1_validation.csv.
% =====================================================================

\section*{Supplementary Methods}
\addcontentsline{toc}{section}{Supplementary Methods}
\label{smethods}

This section gives the procedures that the main text and the
Supplementary Notes refer to, in the order of the main Methods.

\subsection{Severity}
\label{smeth:severity}

Two studies in the corpus record how severe each patient's disease is, and
they record it in different vocabularies. In \texttt{marecko\_2024},
severity is the stage of the revised American Society for Reproductive
Medicine (rASRM) classification \citep{asrm1997revised}, which scores the
number, depth and adhesion of lesions found at surgery into stages I to
IV, with controls coded 0. In \texttt{duempelmann\_2026}, each patient is
graded mild or severe and controls are recorded separately. No other study
in the corpus records severity.

We placed both vocabularies on one three-level scale. Controls from either
study were assigned to \emph{control}. Stages I and II, which the rASRM
classification names minimal and mild disease, were assigned to
\emph{mild} together with the mild grade, and stages III and IV, moderate
and severe disease, to \emph{severe} together with the severe grade. Each
cell takes its value from one source field, because each study supplies
only one of the two, so no cell required the vocabularies to be reconciled
against each other. The merged scale covers 641{,}024 cells from 118
donors: 61 controls, 29 patients with mild disease and 28 with severe
disease.

A contrast that includes controls is partly a test of disease status,
which the latent separates (main text). To ask whether the latent
separates stages of disease among patients, we therefore also used the
same scale with controls removed, which leaves 333{,}984 cells from 57
patients. In addition, the rASRM stage and the grade were each scored
separately within their own study. A single implementation of the mapping
is shared by every analysis and figure in the paper, and all three forms
are given per donor in Supplementary Data~3. Overall, severity is assessed
in three forms: merged with controls, merged among patients only, and in
each study's own vocabulary.

\subsection{Pathway gene grouping}
\label{smeth:grouping}

The model reads genes in groups taken from two curated hierarchies of
biological processes. Gene membership came from the Reactome 2022
\citep{gillespie2022reactome} and Gene Ontology (GO) Biological Process
2023 \citep{go2023knowledgebase} libraries distributed by Enrichr
\citep{kuleshov2016enrichr}, and the parent--child structure from
Reactome's pathway relation file, restricted to human pathways, and from
the GO ontology file. The universe to be grouped was the 30{,}962 genes
found in at least one study of the corpus after symbol resolution (main
Methods, \emph{Gene vocabulary}).

\textbf{One cut, applied twice.} We cut each hierarchy with the same
procedure, which holds every group between a floor of 250 and a ceiling of
500 genes. It begins with the roots of the hierarchy that carry at least
one gene of the universe. First, any node with more than 500 genes is
replaced by its children, repeatedly, until no node exceeds the ceiling.
When a node is replaced, the genes annotated to it directly and to none of
its children would otherwise be lost, so they are kept as a residual
group of that node. Second, every gene is assigned to the most specific
node that contains it. A gene reachable from several nodes, which is
common because both hierarchies allow several parents, goes to the
smallest of them, with ties broken by node identifier so that the result
is the same on every run. Third, groups below the floor are merged. The
smallest group is merged into the partner that shares its most specific
common ancestor, so that a merged group stays one coherent branch of the
hierarchy and can be named after that ancestor. Among equally related
partners the smaller is preferred, and a partner whose merge would exceed
the ceiling is not considered at all. This step repeats until every group
clears the floor or cannot be merged.

The ceiling takes precedence over the floor. It bounds the memory of
attention, which grows with the square of group size, and because groups
are run in size buckets padded to their largest member, one oversized
group would be paid for by every group in its bucket. A group whose every
merge would exceed 500 genes is therefore left below the floor. This
happened once, to a GO group of 247 genes.

\textbf{Reactome first, GO on the residual.} We cut Reactome over the
whole universe first, which assigned 10{,}114 genes to 28 groups. GO was
then cut over the genes Reactome had not assigned, which added 5{,}382
genes in 16 groups. Reactome takes precedence because it is manually curated and organised
by pathway, whereas GO is broader and its annotation noisier.
The 15{,}466 genes annotated in neither hierarchy were dropped. Each group carries the name of one of its constituent pathways, or of an
informative common ancestor, followed by the number of further pathways
it absorbed.

The resulting vocabulary has 15{,}496 genes in 44 groups of 247 to 493
genes, with a median of 342 (Supplementary Table~\ref{stab:vocabulary}).
Every gene belongs to exactly one group, and the full assignment is given
in Supplementary Data~5. In summary, the grouping is fixed before
training, derived only from the two hierarchies, and never adjusted
afterwards.

\subsection{Training configuration}
\label{smeth:hyperparameters}

Supplementary Table~\ref{stab:hyperparameters} lists every setting of the
model reported in the main text. The table is generated from the
configuration that the training run itself recorded, so it gives the
values actually used. The main Methods give the reasons for the settings
that were chosen deliberately: the gene grouping, the bucketed encoder,
the latent width and the weights on each objective. The remaining settings
are conventional defaults that were not tuned.

The checkpoint was selected by reconstruction loss on the held-out
validation donors. Labels played no part in selection, and every metric in
the paper is computed on that one checkpoint, at epoch 4 of a run that
stopped early at epoch 9. Every other configuration in the paper differs
from this one only in the settings named where it is introduced
(\ref{note:frontier} and \ref{note:ablations}).

\begin{table}[t]
\centering
\caption[Training configuration]{\textbf{Every setting of the reported
model.} Generated from the training run's own recorded configuration and
log. KL, Kullback--Leibler divergence. The adversary's gradient is
reversed and scaled by $\lambda_{\mathrm{dab}}$ on its way to the encoder
(main Methods, \emph{Adversarial study removal head}).}
\label{stab:hyperparameters}
\small
% GENERATED by figures/scripts/supp_methods.py from run_config.json. Do not edit.
\begin{tabular}{@{}ll@{}}
\toprule
Setting & Value \\
\midrule
\multicolumn{2}{@{}l}{\textit{Input}} \\
\quad Count scaling & log, library size scaled to 10{,}000 \\
\quad Value encoding & one learned direction per gene, scaled by its count \\
\quad Token width $d$ & 32 \\
\quad Genes, groups & 15{,}496 genes in 44 groups \\
\addlinespace
\multicolumn{2}{@{}l}{\textit{Gene encoder}} \\
\quad Transformer layers per group & 2 \\
\quad Attention heads & 4 \\
\quad Dropout & 0.1 \\
\quad Summary tokens per group & 9 (1 learned, 8 pooled statistics) \\
\quad Size buckets & 4 \\
\addlinespace
\multicolumn{2}{@{}l}{\textit{Cell encoder}} \\
\quad Transformer layers & 2 \\
\quad Attention heads & 4 \\
\quad Dropout & 0.1 \\
\quad Width $d_{\mathrm{agg}}$ & 64 \\
\addlinespace
\multicolumn{2}{@{}l}{\textit{Latent}} \\
\quad Gaussian dimensions & 16 \\
\quad Circular dimensions & 0 \\
\addlinespace
\multicolumn{2}{@{}l}{\textit{Decoder}} \\
\quad Width $d_{\mathrm{dec}}$ & 64 \\
\quad Tissue, platform, study embeddings & 16, 16, 16 dimensions \\
\quad Covariates zeroed per step & probability 0.25 \\
\addlinespace
\multicolumn{2}{@{}l}{\textit{Objective}} \\
\quad Reconstruction weight $w_{\mathrm{nb}}$ & 1 \\
\quad KL weight $w_{\mathrm{kl}}$ & 1 \\
\quad KL ramp & 0 to 0.5 over the first 5\% of steps \\
\quad KL free bits & 0.25 \\
\quad Adversary weight $w_{\mathrm{dab}}$, reversal $\lambda_{\mathrm{dab}}$ & 0.5, 168 \\
\quad Adversary network & 3 layers, 256 hidden units \\
\quad Phase head weight, label smoothing & 40, 0.05 \\
\quad Cell-similarity weight, threshold & 50, 0.8 \\
\addlinespace
\multicolumn{2}{@{}l}{\textit{Optimisation}} \\
\quad Optimiser & AdamW, $\beta = (0.9, 0.999)$ \\
\quad Learning rate, weight decay & $10^{-3}$, $10^{-4}$ \\
\quad Schedule & cosine over 20 epochs \\
\quad Batch size & 256 \\
\quad Gradient clipping & norm 1 \\
\quad Precision & bf16 \\
\quad Early stopping & held-out reconstruction, patience 5 after 2 epochs \\
\quad Selected epoch & 4 (held-out reconstruction 0.4024) \\
\addlinespace
\multicolumn{2}{@{}l}{\textit{Data}} \\
\quad Split & by donor, 70/15/15, seed 24 \\
\quad Training seed & 24 \\
\quad Parameters & 1{,}883{,}986 \\
\bottomrule
\end{tabular}

\end{table}

\subsection{Label transfer}
\label{smeth:transfer}

Two studies annotated their own cells. \texttt{duempelmann\_2026} and
\texttt{marecko\_2024} supply a cell family for 640{,}606 cells from 118
donors, while the other nine studies supply none. To compare cell families
across the whole corpus, we assigned a family to each of the remaining
937{,}122 cells by nearest-neighbour vote in the 16-dimensional latent of
the reported model. The transfer is restricted to the four families
stromal, epithelial, endothelial and immune. The two studies annotate
finer cell types in vocabularies of 45 and 63 types that do not agree
with each other, and nothing in the corpus could validate a transfer at
that resolution.

\textbf{Procedure.} For each unannotated cell, we found its 25 nearest
annotated cells by Euclidean distance in the latent and assigned the
family held by most of them, with ties broken in a fixed order. Every
transfer is therefore across both donors and studies, since the nine
target studies contribute no annotated cells. For each assignment we also
recorded the share of the 25 neighbours that voted for the winning family,
their mean distance, and how many donors and studies they came from, so
that any downstream analysis can restrict itself to confident assignments.
Annotated cells keep their own label. The search used scikit-learn
\citep{pedregosa2011scikit}.

\textbf{Validation before transfer.} Two tests had to pass before any label
was written. First, we transferred labels between the two annotated
studies in each direction, using 60{,}000 cells of one study as the
reference and 60{,}000 of the other as the test set, and required a
balanced accuracy above 0.90. Transfer from \texttt{marecko\_2024} to
\texttt{duempelmann\_2026} reached an accuracy of 0.9862 and a balanced
accuracy of 0.9596, against a majority baseline of 0.8805. The reverse
direction reached 0.9845 and 0.9640, against 0.5252. Second, we
transferred labels into \texttt{li\_2025}, a follicular-fluid study whose
own annotation uses a different vocabulary, and required its immune cells
to be recovered as immune. They were, at 0.9966. The same test also showed
the procedure's failure mode. Granulosa cells, which have no counterpart in
endometrium, were absorbed into an existing family and were not flagged as
unfamiliar (\ref{note:limitations}). Overall, the
transfer recovers held-out annotations across studies with high accuracy,
and a cell type absent from the annotated studies is assigned to the
nearest family that is present.
